%!TEX root = main.tex
% ============================================================
% K12 教辅绘图项目 — 公共导言区
% 编译方式：xelatex main.tex
% ============================================================

\documentclass[12pt, a4paper]{ctexrep}
\xeCJKsetup{CJKmath=true}
\xeCJKDeclareCharClass{CJK}{"2460, "2461, "2462, "2463, "2464, "2465, "25B1, "3001}

\usepackage{geometry}
\geometry{
  a4paper,
  left=2.5cm,
  right=2.5cm,
  top=2.5cm,
  bottom=2.5cm
}

\usepackage{amsmath}
\usepackage{amssymb}
\usepackage{array}
\usepackage{booktabs}
\usepackage{multirow}
\usepackage{makecell}
\usepackage{colortbl}
\usepackage{hhline}
\usepackage{float}
\usepackage{tikz}
\usetikzlibrary{
  angles,                   % 画角与直角标记
  quotes,                   % 给角附文字标签
  arrows.meta,              % Stealth 等箭头样式
  calc,                     % 坐标计算
  decorations.pathreplacing,% 大括号装饰
  decorations.markings,     % 线段标记装饰（附录参考绘图）
  decorations.pathmorphing, % 波浪线等路径变形（附录参考绘图）
  patterns,                 % 阴影填充图案
  positioning,              % 节点相对定位（附录参考绘图）
  snakes                    % 弹簧、波浪装饰（附录参考绘图，兼容旧版）
}

% ------ 附录参考绘图所需额外宏包 ------
\usepackage{tkz-euclide}          % 欧几里得几何作图（jamesfang8499 数学教材）
\usepackage{circuitikz}           % 电路图绘制（jamesfang8499 物理教材）
\usepackage{graphicx}             % 图片引用
\usepackage{longtable}            % 跨页跨列表格
\usepackage{multirow}             % 多列表格支持
\usepackage{wasysym}              % 为其它特殊符号提供支持
\usepackage{marvosym}             % 提供 \Sun 等符号支持
\usepackage[normalem]{ulem}       % 为 \uwave 波浪下划线提供支持

% ------ 原教材缺失的自定义宏 ------
\newcommand{\x}{\times}
\newcommand{\ms}{\mathrm{m/s}}
\newcommand{\msq}{\mathrm{m/s}^2}

% ------ 正字计数法自定义符号 ------
\newcommand{\zhengI}{\tikz[baseline=0.1em, x=1em, y=1em]{\draw[line width=1pt, line cap=round] (0.15, 0.8) -- (0.85, 0.8);}}
\newcommand{\zhengII}{\tikz[baseline=0.1em, x=1em, y=1em]{\draw[line width=1pt, line cap=round] (0.15, 0.8) -- (0.85, 0.8); \draw[line width=1pt, line cap=round] (0.5, 0.8) -- (0.5, 0.1);}}
\newcommand{\zhengIII}{\tikz[baseline=0.1em, x=1em, y=1em]{\draw[line width=1pt, line cap=round] (0.15, 0.8) -- (0.85, 0.8); \draw[line width=1pt, line cap=round] (0.5, 0.8) -- (0.5, 0.1); \draw[line width=1pt, line cap=round] (0.5, 0.45) -- (0.8, 0.45);}}
\newcommand{\zhengIV}{\tikz[baseline=0.1em, x=1em, y=1em]{\draw[line width=1pt, line cap=round] (0.15, 0.8) -- (0.85, 0.8); \draw[line width=1pt, line cap=round] (0.5, 0.8) -- (0.5, 0.1); \draw[line width=1pt, line cap=round] (0.5, 0.45) -- (0.8, 0.45); \draw[line width=1pt, line cap=round] (0.25, 0.45) -- (0.25, 0.1);}}
\newcommand{\zhengV}{\tikz[baseline=0.1em, x=1em, y=1em]{\draw[line width=1pt, line cap=round] (0.15, 0.8) -- (0.85, 0.8); \draw[line width=1pt, line cap=round] (0.5, 0.8) -- (0.5, 0.1); \draw[line width=1pt, line cap=round] (0.5, 0.45) -- (0.8, 0.45); \draw[line width=1pt, line cap=round] (0.25, 0.45) -- (0.25, 0.1); \draw[line width=1pt, line cap=round] (0.05, 0.1) -- (0.95, 0.1);}}
\newcommand{\zheng}{\zhengV}

% ------ TikZ 全局样式 ------
\tikzset{
  ray/.style={
    -{Stealth[length=6pt]},
    line width=1.2pt,
    color=cyan!70!blue
  },
  given ray/.style={
    -,
    line width=1.2pt,
    color=cyan!70!blue
  },
  angle arc/.style={
    draw=red!70!black,
    line width=0.8pt
  },
  angle label/.style={
    font=\small,
    color=red!70!black
  },
  vertex dot/.style={
    circle,
    fill=cyan!70!blue,
    inner sep=0pt,
    minimum size=5pt
  }
}

%% ------ 实心五角星宏（统计图用） ------
\newcommand{\tikzSolidStar}[1]{%
  \tikz[baseline=-0.3em]{%
    \node[#1, inner sep=0pt, font=\large] at (0,0.1) {$\bigstar$};%
  }%
}

% ------ 抑制烦人的弱警告 ------
\usepackage{silence}
\WarningFilter{hyperref}{Token not allowed in a PDF string}

% ------ 超链接设置 ------
\usepackage[colorlinks=true, linkcolor=blue, filecolor=blue, urlcolor=blue, citecolor=blue]{hyperref}
